\section{Discussion}
\label{sec:discussion}

The limitations below are pre-committed: they were written before any
result was read, so that none can be selected after the fact to excuse a
finding or to inflate one.

\subsection{The effective sample size is one}

Every shift-regime number in this paper is a statement about one pandemic
observed in twenty places. Clustering by population and imposing the null
in the bootstrap make the standard errors honest \emph{given the event}; they
do not make the event a sample of events. The audit therefore answers
``did these intervals survive this break?'' and, through the placebo,
``does the pattern resemble the previous break?''---not ``will these
intervals survive the next break?''. A reader who wants the last question
answered wants a different study, one that we cannot run until there is a
next break.

\subsection{The grid is ragged, so claims are contrasts}

Ten families by seven mechanisms is seventy cells on paper and fifty in
fact. Main effects over a ragged grid confound the factor with the pattern
of admissibility, so no main effect is reported; every claim is a contrast
within a sub-grid whose composition is stated. This costs generality: a
finding that conformal wrapping changes coverage for classical families
says nothing about neural families unless the neural sub-grid shows it
too, and we will not average across the two.

\subsection{Conformal cells carry secondary proper scores}

A conformal band is an interval, not a distribution. Drawing uniformly
inside it keeps the single interface honest for coverage, width and
interval score, but the CRPS, log score and PIT of those draws describe a
placeholder density. They are reported in Appendix~\ref{app:secondary},
flagged, and not ranked. Any reader who compares a conformal cell's CRPS
with a native cell's CRPS is comparing a convention with a forecast.
Conformal predictive systems that output a full calibrated distribution
would remove this limitation and are noted as future work.

\subsection{2024 may be revised}

HMD does not flag provisional years; it re-issues whole series at each
release, so the 2024 values are final only in the sense that this release
contains them. The shift regime's $h=5$ results rest on 2024, and two of the
twenty populations (USA, CHL) carry census-based exposures that HMD
documents as pending revision. The pre-registered robustness run drops 2024;
Addendum~1 adds runs that drop USA and CHL and that contrast register-based
with census-based populations. If the $h=5$ and $h\le4$ results disagree,
the $h\le4$ result is the one to trust and the $h=5$ result is reported as
provisional; if a later vintage appears before submission, the shift regime
is re-run on it and any change in the headline coverage is reported as a
vintage-revision effect.

\subsection{SVD Lee--Carter is over-dispersed by construction}

Validation gate~1 found that the classical two-stage SVD Lee--Carter
produces intervals slightly wider than nominal on data simulated from a
Poisson Lee--Carter process, because the SVD estimate of $\kappa_t$ absorbs
observation noise that then inflates the random-walk innovation variance.
This is a property of the method that \citet{girosi2007understanding}
would predict, not a harness artefact, and it has a direct consequence for
reading the results: an SVD-LC row that appears better calibrated than
Poisson-LC in the stable regime may be so partly because it is vaguer, and
the interval score and Murphy resolution term---not coverage---are the
columns that tell the two apart.

\subsection{Other limitations}

\begin{itemize}
\item \emph{Exchangeability is violated by design.} Conformal calibration
  residuals come from pre-cutoff years, and the test years are precisely
  the years chosen to differ from them. \citet{barber2023conformal} bound the
  resulting coverage gap; we measure it rather than assume it away. The
  conformal arms are audited, not endorsed.
\item \emph{Neural families are budget-limited.} Ensembles of ten over
  thirteen origins, twenty populations and two sexes fix the hyperparameter
  search to what fits the GPU budget; the inner time split is small (five
  years). A better-tuned network might do better on point accuracy; whether
  it would be better calibrated is exactly what is not known.
\item \emph{Data quality in the placebo window.} Military deaths are
  reconstructed in the belligerent series and excluded in the civilian-only
  one. Placebo results are descriptive, are reported pooled and by the
  pre-specified strata, and the neutral register-based stratum carries the
  1918-versus-2020 comparison.
\item \emph{Single-age granularity.} Modelling ages $0$--$99$ singly keeps
  the old-age cells noisy in small populations; the age-cap-90 robustness
  run tests whether \Hfour\ is an artefact of that noise.
\item \emph{Rounding of deaths.} The log-score convention rounds fractional
  deaths; the CRPS-on-counts sensitivity exists to show that nothing turns
  on it.
\item \emph{One code path, one author.} A single evaluation module is a
  strength for consistency and a weakness for independent verification; the
  synthetic-truth gate and the \pkg{StMoMo} oracle are the checks, and the
  code is public.
\end{itemize}

\subsection{What this paper does not do}

It does not propose a model, a repair, or a recommendation about which
family to use. If the audit finds that a mechanism restores coverage at the
break, that is a finding about the mechanism to be stated as a contrast; it
is not a method contribution and will not be presented as one. If the audit
finds that nothing holds nominal coverage, that is the finding, and the
actuarial reading in Section~\ref{sec:actuarial} is its consequence.
Explainability of the neural families---which features drive their
forecasts and whether the explanations are stable across the break---is a
separate question deferred to a companion paper that reuses these fits.
